\documentclass[12pt,a4paper]{article}
\usepackage[utf8]{vietnam}
\usepackage{amsmath,amssymb}
\usepackage{geometry}
\geometry{top=1.3cm,bottom=1.2cm,left=1.4cm,right=1.4cm,headheight=15pt,headsep=12pt,footskip=18pt}
\usepackage{tikz}
\usepackage{xcolor}
\usepackage{enumerate}
\usepackage{graphicx}
\usepackage[version=4]{mhchem}
\usepackage{fancyhdr}
\usepackage{ex_test}

\renewcommand{\baselinestretch}{0.95}
\setlength{\parskip}{0pt}

\definecolor{hfRed}{RGB}{211,47,47}
\definecolor{hfGreen}{RGB}{139,195,144}

\pagestyle{fancy}
\fancyhf{}
\fancyhead[L]{\small\itshape\color{hfRed} Tài liệu lưu hành nội bộ}
\fancyhead[R]{\small\itshape\color{hfRed} Thầy \textbf{TRẦN VĂN THẠNH} - 0777.470.803}
\renewcommand{\headrulewidth}{0.8pt}
\renewcommand{\headrule}{\hbox to\headwidth{\color{hfGreen}\leaders\hrule height \headrulewidth\hfill}}

\fancyfoot[L]{\small\itshape\color{hfRed} CS1: 6/15 Nguyễn Hoàng, P. Kim Long, TP Huế \quad|\quad CS2: 24 Đặng Thái Thân, TP Huế}
\fancyfoot[R]{\small\bfseries\color{hfRed} Trang \thepage}
\renewcommand{\footrulewidth}{0.8pt}
\renewcommand{\footrule}{\hbox to\headwidth{\color{hfGreen}\leaders\hrule height \footrulewidth\hfill}}

\fancypagestyle{firstpage}{
  \fancyhf{}
  \renewcommand{\headrulewidth}{0pt}
  \fancyfoot[L]{\small\itshape\color{hfRed} CS1: 6/15 Nguyễn Hoàng, P. Kim Long, TP Huế \quad|\quad CS2: 24 Đặng Thái Thân, TP Huế}
  \fancyfoot[R]{\small\bfseries\color{hfRed} Trang \thepage}
  \renewcommand{\footrulewidth}{0.8pt}
  \renewcommand{\footrule}{\hbox to\headwidth{\color{hfGreen}\leaders\hrule height \footrulewidth\hfill}}
}

\begin{document}
\thispagestyle{firstpage}

\noindent
\begin{minipage}[t]{0.42\textwidth}
	\centering\fontsize{9.5pt}{12pt}\selectfont
	\textbf{SỞ GD \& ĐT THÀNH PHỐ HUẾ}\\
	\textbf{TRƯỜNG THPT NGUYỄN CHÍ THANH}\\
	\textbf{ĐỀ CHÍNH THỨC}\\
	\textit{(Đề thi có 04 trang)}
\end{minipage}\hfill
\begin{minipage}[t]{0.56\textwidth}
	\centering\small
	\textbf{ĐỀ KIỂM TRA GIỮA HỌC KÌ 1 - NĂM HỌC 2025--2026}\\
	\textbf{MÔN: HÓA HỌC -- 12CTST}\\
	\textit{Thời gian làm bài: 45 phút (không kể thời gian phát đề)}
\end{minipage}

\smallskip
\noindent\hrulefill
\smallskip

\noindent\textbf{Họ và tên:} \dotfill\ \textbf{Số báo danh:} \dotfill\ \textbf{Mã đề 104}

\smallskip
\noindent\textit{Cho biết nguyên tử khối của các nguyên tố (amu): H=1; O=16; C=12; Ag=108; Na=23; N=14.}
\smallskip

\noindent
\textbf{A. PHẦN TRẮC NGHIỆM (7 điểm)}

\smallskip
\noindent
\textbf{Phần I. Câu trắc nghiệm nhiều phương án lựa chọn.} \textit{(Thí sinh trả lời từ câu 1 đến câu 16. Mỗi câu hỏi thí sinh chỉ chọn một phương án).}

\Opensolutionfile{ans}[ans-gv]

\begin{ex}
Phản ứng nào dưới đây \textbf{không} thể hiện tính base của amine?
\choice
{$\ce{C6H5NH2 + HCl -> C6H5NH3Cl}$}
{\True $\ce{RNH2 + HNO2 -> ROH + N2 ^ + H2O}$}
{$\ce{Fe^{3+} + 3RNH2 + 3H2O -> Fe(OH)3 v + 3RNH3^+}$}
{$\ce{RNH2 + H2O <=> RNH3^+ + OH^-}$}
\loigiai{
Phản ứng của amine bậc một tác dụng với acid nitrous $\ce{HNO2}$ sinh ra alcohol $\ce{ROH}$ và giải phóng khí $\ce{N2}$ là phản ứng thế nhóm amino bằng nhóm -OH, không thể hiện tính base của amine.\\
\textbf{Chọn B.}
}
\end{ex}

\begin{ex}
Nhận xét nào dưới đây là không đúng khi nói về glucose và fructose?
\choice
{Đều tạo được kết tủa đỏ gạch $\ce{Cu2O}$ khi tác dụng với $\ce{Cu(OH)2}$, đun nóng trong môi trường kiềm}
{Đều xảy ra phản ứng tráng bạc khi tác dụng với thuốc thử Tollens}
{Đều tạo được dung dịch màu xanh lam khi tác dụng với $\ce{Cu(OH)2}$ trong môi trường kiềm}
{\True Đều làm mất màu nước bromine}
\loigiai{
Fructose không làm mất màu nước bromine vì fructose không chứa nhóm aldehyde (-CHO), và trong môi trường nước bromine có tính acid yếu nên không có sự chuyển hóa giữa fructose và glucose.\\
\textbf{Chọn D.}
}
\end{ex}

\begin{ex}
Khi đun nóng chất X có công thức phân tử $\ce{C3H6O2}$ với dung dịch $\ce{NaOH}$ thu được $\ce{CH3COONa}$. Công thức cấu tạo của X là:
\choice
{$\ce{CH3COOC2H5}$}
{\True $\ce{CH3COOCH3}$}
{$\ce{HCOOC2H5}$}
{$\ce{C2H5COOH}$}
\loigiai{
Muối sinh ra là $\ce{CH3COONa}$ (chứa 2 nguyên tử C). Do X có 3 carbon nên gốc alcohol còn lại có 1 carbon ($\ce{-CH3}$). Vậy CTCT của X là $\ce{CH3COOCH3}$ (methyl acetate).\\
\textbf{Chọn B.}
}
\end{ex}

\begin{ex}
Trong cây thuốc lá tự nhiên và khói thuốc lá chứa một amine rất độc, đó là nicotin với công thức cấu tạo như sau:
\begin{center}
	\includegraphics[width=13.5cm]{San_Pham/images_crop/fig_cau4_nicotin.png}
\end{center}
Nicotin làm tăng huyết áp và nhịp tim, có khả năng gây xơ vữa động mạch vành và suy giảm trí nhớ. Số nguyên tử carbon trong một phân tử nicotin là:
\choice
{11}
{8}
{9}
{\True 10}
\loigiai{
Vòng pyridine gồm 5 nguyên tử C; vòng pyrrolidine gồm 4 nguyên tử C trong vòng và 1 nguyên tử C ở nhóm methyl ($\ce{-CH3}$) gắn vào nguyên tử N. Tổng số nguyên tử carbon trong phân tử nicotin = $5 + 4 + 1 = 10$ (CTPT là $\ce{C10H14N2}$).\\
\textbf{Chọn D.}
}
\end{ex}

\begin{ex}
Chất giặt rửa tổng hợp sodium laurylsulfate có công thức cấu tạo như sau:
\begin{center}
	\includegraphics[width=9.5cm]{San_Pham/images_crop/fig_cau5_sodium_laurylsulfate.png}
\end{center}
Nhóm được khoanh tròn trong công thức trên là
\choice
{đầu kị nước}
{đuôi ưa nước}
{\True đầu ưa nước}
{đuôi kị nước}
\loigiai{
Phân tử sodium laurylsulfate gồm đuôi kị nước (mạch hydrocarbon dài $\ce{CH3[CH2]10CH2-}$) và đầu ưa nước là nhóm ion sulfate phân cực $\ce{-OSO3^-Na^+}$ (được khoanh tròn trong hình).\\
\textbf{Chọn C.}
}
\end{ex}

\begin{ex}
Chất nào sau đây là amine bậc hai?
\choice
{$\ce{(C2H5)3N}$}
{$\ce{CH3CH(NH2)CH3}$}
{$\ce{C2H5NH2}$}
{\True $\ce{(C2H5)2NH}$}
\loigiai{
Amine bậc hai có 2 gốc hydrocarbon liên kết với nguyên tử nitrogen: $\ce{(C2H5)2NH}$ (diethylamine).\\
\textbf{Chọn D.}
}
\end{ex}

\begin{ex}
Bệnh nhân phải tiếp đường (truyền dung dịch đường vào tĩnh mạch), đó là loại đường nào?
\choice
{\True Glucose}
{Cellulose}
{Saccharose}
{Fructose}
\loigiai{
Dung dịch glucose 5\% (huyết thanh ngọt) được dùng để truyền trực tiếp vào tĩnh mạch cung cấp năng lượng nhanh cho cơ thể suy nhược.\\
\textbf{Chọn A.}
}
\end{ex}

\begin{ex}
$\ce{HCOOCH3}$ được sử dụng làm dung môi cho nitrocellulose và cellulose acetate, một chất trung gian trong sản xuất dược phẩm và thuốc xông hơi. Tên gọi của $\ce{HCOOCH3}$ là:
\choice
{methyl acetate}
{\True methyl formate}
{ethyl acetate}
{ethyl formate}
\loigiai{
$\ce{HCOOCH3}$ có gốc alcohol là methyl ($\ce{-CH3}$) và gốc acid là formate ($\ce{HCOO-}$) $\rightarrow$ Tên gọi là methyl formate.\\
\textbf{Chọn B.}
}
\end{ex}

\begin{ex}
Xà phòng hoá hoàn toàn $17{,}24\text{ gam}$ chất béo cần vừa đủ $0{,}06\text{ mol}$ $\ce{NaOH}$. Cô cạn dung dịch sau phản ứng thu được bao nhiêu gam xà phòng?
\choice
{$21{,}15\text{ gam}$}
{$14{,}12\text{ gam}$}
{$14{,}84\text{ gam}$}
{\True $17{,}8\text{ gam}$}
\loigiai{
Phương trình xà phòng hóa: Chất béo + $3\ce{NaOH} \rightarrow$ Xà phòng + Glycerol\\
$n_{\text{glycerol}} = \dfrac{1}{3} n_{\ce{NaOH}} = \dfrac{0{,}06}{3} = 0{,}02\text{ mol}$.\\
Áp dụng định luật bảo toàn khối lượng:
\[m_{\text{chất béo}} + m_{\ce{NaOH}} = m_{\text{xà phòng}} + m_{\text{glycerol}}\]
\[\Rightarrow m_{\text{xà phòng}} = 17{,}24 + 0{,}06 \times 40 - 0{,}02 \times 92 = 17{,}24 + 2{,}40 - 1{,}84 = 17{,}80\text{ gam}.\]
\textbf{Chọn D.}
}
\end{ex}

\begin{ex}
Chất nào sau đây là thành phần chính của xà phòng?
\choice
{$\ce{CH3[CH2]_{16}COOCH3}$}
{\True $\ce{CH3[CH2]_{14}COONa}$}
{$\ce{CH3COONa}$}
{$\ce{CH3[CH2]_{11}C6H4SO3Na}$}
\loigiai{
Thành phần chính của xà phòng là muối sodium hoặc potassium của các acid béo. $\ce{CH3[CH2]14COONa}$ là sodium palmitate.\\
\textbf{Chọn B.}
}
\end{ex}

\begin{ex}
Aniline tác dụng với $(\ce{HNO2 + HCl})$ ở $0 - 5^\circ\text{C}$ tạo muối diazonium để tổng hợp phẩm nhuộm azo và dược phẩm.
\[\ce{C6H5NH2 + HONO + HCl ->[0 - 5^\circ\text{C}] X + 2H2O}\]
Chất X có công thức cấu tạo là:
\choice
{$[\ce{C6H5N2H}]^+\ce{Cl^-}$}
{$[\ce{C6H5NH2}]^+\ce{Cl^-}$}
{\True $[\ce{C6H5N2}]^+\ce{Cl^-}$}
{$[\ce{C6H5NH3}]^+\ce{Cl^-}$}
\loigiai{
Phản ứng diazot hóa amine thơm bậc một ở $0 - 5^\circ\text{C}$ tạo muối benzenediazonium chloride: $[\ce{C6H5N2}]^+\ce{Cl^-}$.\\
\textbf{Chọn C.}
}
\end{ex}

\begin{ex}
Các gốc $\alpha$-glucose trong phân tử tinh bột tạo dạng mạch amylopectin phân nhánh, xoắn. Phần phân nhánh liên kết với nhau bởi liên kết
\choice
{$\beta$-1,2-glycoside}
{$\alpha$-1,4-glycoside}
{$\alpha$-1,3-glycoside}
{\True $\alpha$-1,6-glycoside}
\loigiai{
Trong phân tử amylopectin, các gốc $\alpha$-glucose trong mạch liên kết bằng liên kết $\alpha$-1,4-glycoside; tại các điểm phân nhánh, mạch nhánh liên kết với mạch chính bằng liên kết $\alpha$-1,6-glycoside.\\
\textbf{Chọn D.}
}
\end{ex}

\begin{ex}
Một số ester được dùng trong hương liệu, mĩ phẩm, bột giặt là nhờ các ester
\choice
{\True có mùi thơm, an toàn với người sử dụng}
{đều có nguồn gốc từ thiên nhiên}
{có thể bay hơi nhanh sau khi sử dụng}
{là chất lỏng dễ bay hơi}
\loigiai{
Nhiều ester có mùi thơm hoa quả đặc trưng, dễ chịu và an toàn cho người sử dụng nên được dùng làm hương liệu cho mĩ phẩm, thực phẩm, bột giặt.\\
\textbf{Chọn A.}
}
\end{ex}

\begin{ex}
Cho dung dịch ethylamine lần lượt tác dụng với: dung dịch $\ce{FeCl3}$; dung dịch $\ce{HCl}$; $\ce{Cu(OH)2}$; dung dịch $\ce{NaOH}$. Số trường hợp xảy ra phản ứng là:
\choice
{2}
{4}
{5}
{\True 3}
\loigiai{
Ethylamine $\ce{C2H5NH2}$ có tính base, phản ứng được với 3 chất:
\begin{enumerate}[(1)]
	\item Dung dịch $\ce{FeCl3}$: tạo kết tủa nâu đỏ $\ce{Fe(OH)3}$.
	\item Dung dịch $\ce{HCl}$: tạo muối $\ce{C2H5NH3Cl}$.
	\item $\ce{Cu(OH)2}$: hòa tan tạo phức đồng có màu xanh lam thẫm $[\ce{Cu(C2H5NH2)4}](\ce{OH})_2$.
\end{enumerate}
Không phản ứng với dung dịch $\ce{NaOH}$. Vậy có 3 trường hợp xảy ra phản ứng.\\
\textbf{Chọn D.}
}
\end{ex}

\begin{ex}
Kết quả thí nghiệm của các dung dịch X, Y, Z với thuốc thử được ghi ở bảng sau:
\begin{center}
\begin{tabular}{|c|l|l|}
\hline
\textbf{Mẫu thử} & \textbf{Thuốc thử} & \textbf{Hiện tượng} \\ \hline
X & Dung dịch $\ce{I2}$ & Có màu xanh tím \\ \hline
Y & $\ce{Cu(OH)2}$ & Có màu xanh lam \\ \hline
Z & Dung dịch $\ce{AgNO3}$ trong $\ce{NH3}$ & Tạo kết tủa Ag \\ \hline
\end{tabular}
\end{center}
Các dung dịch X, Y, Z lần lượt là
\choice
{\True Hồ tinh bột, saccharose, glucose}
{Cellulose, saccharose, glucose}
{Cellulose, glucose, saccharose}
{Hồ tinh bột, glucose, saccharose}
\loigiai{
- X cho màu xanh tím với dung dịch $\ce{I2} \rightarrow$ X là Hồ tinh bột.\\
- Y hòa tan $\ce{Cu(OH)2}$ ở nhiệt độ thường cho dung dịch màu xanh lam $\rightarrow$ Y là Saccharose.\\
- Z phản ứng với dung dịch $\ce{AgNO3}$ trong $\ce{NH3}$ tạo kết tủa $\ce{Ag} \rightarrow$ Z là Glucose.\\
\textbf{Chọn A.}
}
\end{ex}

\begin{ex}
Công thức của triolein là:
\choice
{$\ce{(C17H35COO)3C3H5}$}
{\True $\ce{(C17H33COO)3C3H5}$}
{$\ce{(CH3COO)3C3H5}$}
{$\ce{(HCOO)3C3H5}$}
\loigiai{
Triolein là triglyceride của oleic acid ($\ce{C17H33COOH}$), công thức phân tử thu gọn là $\ce{(C17H33COO)3C3H5}$.\\
\textbf{Chọn B.}
}
\end{ex}

\Closesolutionfile{ans}

\vspace{0.3cm}
\noindent
\textbf{Phần II. Câu trắc nghiệm đúng/sai. (2 điểm -- 2 câu).} \textit{Thí sinh trả lời từ câu 1 đến câu 2. Trong mỗi ý a), b), c), d) ở mỗi câu, thí sinh chọn đúng hoặc sai.}

\setcounter{ex}{0}

\begin{ex}
Tinh bột là loại lương thực quan trọng và là nguyên liệu chủ yếu để sản xuất bánh, kẹo, rượu, bia, ... Cellulose là thành phần chính tạo nên màng tế bào thực vật, có nhiều trong gỗ, bông gòn dùng làm nguyên liệu sản xuất thuốc súng không khói.
\begin{enumerate}[a)]
	\item Tinh bột và cellulose là đồng phân cấu tạo của nhau do cùng có công thức phân tử dạng $\ce{(C6H10O5)_n}$.
	\item Tinh bột và cellulose đều thuộc loại polysaccharide.
	\item Dung dịch hồ tinh bột tạo với iodine hợp chất màu xanh tím, cellulose không có tính chất này.
	\item Thuỷ phân hoàn toàn 162 gam tinh bột hoặc cellulose đều thu được 180 gam sản phẩm là glucose.
\end{enumerate}
\loigiai{
\textbf{Lời giải:}\\
a) \textbf{Sai.} Do hệ số polimer hóa $n$ của tinh bột và cellulose khác nhau rất nhiều ($n_{\text{cellulose}} \gg n_{\text{tinh bột}}$), nên công thức phân tử khác nhau, chúng không phải là đồng phân của nhau.\\
b) \textbf{Đúng.} Cả tinh bột và cellulose đều là polysaccharide được cấu tạo từ nhiều gốc monosaccharide.\\
c) \textbf{Đúng.} Tinh bột có cấu trúc xoắn tạo các khoang rỗng giữ phân tử iodine cho màu xanh tím; cellulose mạch kéo dài không xoắn nên không có tính chất này.\\
d) \textbf{Đúng.} Phương trình: $\ce{(C6H10O5)_n + n H2O -> n C6H12O6}$. Với 162 gam polysaccharide thì $m_{\text{glucose}} = 162 \times \dfrac{180}{162} = 180\text{ gam}$.
}
\end{ex}

\begin{ex}
Aspirin được sử dụng làm thuốc giảm đau, hạ sốt. Sau khi uống, aspirin bị thuỷ phân trong cơ thể tạo thành salicylic acid. Salicylic acid ức chế quá trình sinh tổng hợp prostaglandin (chất gây đau, sốt và viêm khi nồng độ trong máu cao hơn mức bình thường).
\begin{center}
	\includegraphics[width=11.5cm]{San_Pham/images_crop/fig_cau2_p2_aspirin.png}
\end{center}
\begin{enumerate}[a)]
	\item 1 mol aspirin hay 1 mol salicylic acid đều phản ứng tối đa với 2 mol $\ce{NaOH}$.
	\item Trong phân tử aspirin có số liên kết $\pi$ và vòng là 5.
	\item Aspirin có công thức phân tử $\ce{C9H8O4}$.
	\item Aspirin được điều chế theo phản ứng sau:
	\[\ce{(CH3CO)2O + HOC6H4COOH ->[H2SO4, t^\circ] CH3COOC6H4COOH + CH3COOH}\]
	Để sản xuất 3 triệu viên thuốc aspirin cần tối thiểu 260 kg salicylic acid. Biết rằng mỗi viên thuốc có chứa 81 mg aspirin và hiệu suất phản ứng đạt 70\%.
\end{enumerate}
\loigiai{
\textbf{Lời giải:}\\
a) \textbf{Sai.} Aspirin chứa 1 nhóm -COOH và 1 nhóm ester của phenol nên 1 mol aspirin phản ứng tối đa với 3 mol $\ce{NaOH}$:
\[\ce{CH3COOC6H4COOH + 3NaOH -> CH3COONa + NaOC6H4COONa + 2H2O}\]
Salicylic acid chỉ phản ứng với 2 mol $\ce{NaOH}$.\\
b) \textbf{Sai.} Vòng benzene có 1 vòng và 3 liên kết $\pi$; nhóm -COO- có 1 liên kết $\pi$; nhóm -COOH có 1 liên kết $\pi$. Tổng số liên kết $\pi$ và vòng $= 3 + 1 + 1 + 1 = 6$.\\
c) \textbf{Đúng.} Phân tử aspirin ($\ce{CH3COOC6H4COOH}$) gồm 9 C, 8 H, 4 O $\rightarrow \ce{C9H8O4}$.\\
d) \textbf{Sai.} Khối lượng aspirin trong 3 triệu viên: $m_{\text{aspirin}} = 3 \cdot 10^6 \times 81\text{ mg} = 243\text{ kg}$.\\
Lượng salicylic acid cần theo lý thuyết: $m_{\text{acid (LT)}} = 243 \times \dfrac{138}{180} = 186{,}3\text{ kg}$.\\
Do hiệu suất $H = 70\%$, lượng salicylic acid thực tế cần tối thiểu:
\[m_{\text{acid (TT)}} = \dfrac{186{,}3}{0{,}7} \approx 266{,}14\text{ kg} \text{ (khác 260 kg)}.\]
}
\end{ex}

\vspace{0.3cm}
\noindent
\textbf{Phần III. Câu trắc nghiệm yêu cầu trả lời ngắn. (4 câu -- 1 điểm).} \textit{Thí sinh trả lời từ câu 1 đến câu 4.}

\setcounter{ex}{0}

\begin{ex}
Cho công thức cấu tạo thu gọn và tên gọi sau:
\begin{center}
\begin{tabular}{|c|l|l|}
\hline
\textbf{Chất} & \textbf{Công thức cấu tạo thu gọn} & \textbf{Tên gọi} \\ \hline
(1) & $\ce{CH3N(CH3)2}$ & Aniline \\ \hline
(2) & $\ce{CH3CH2NHCH3}$ & Trimethylamine \\ \hline
(3) & $\ce{CH3CH2NH2}$ & Ethylamine \\ \hline
(4) & $\ce{C6H5NH2}$ & N-methylethanamine \\ \hline
\end{tabular}
\end{center}
Sắp xếp các chất theo thứ tự tên gọi tương ứng?
\loigiai{
Đối chiếu công thức cấu tạo thu gọn với tên gọi đúng:
\begin{itemize}
	\item (1) $\ce{CH3N(CH3)2}$ là Trimethylamine (tên dòng 2).
	\item (2) $\ce{CH3CH2NHCH3}$ là N-methylethanamine (tên dòng 4).
	\item (3) $\ce{CH3CH2NH2}$ là Ethylamine (tên dòng 3).
	\item (4) $\ce{C6H5NH2}$ là Aniline (tên dòng 1).
\end{itemize}
Thứ tự các chất tương ứng với cột tên gọi từ trên xuống dưới (Aniline, Trimethylamine, Ethylamine, N-methylethanamine) là: \textbf{(4) - (1) - (3) - (2)}.
}
\end{ex}

\begin{ex}
Tại một nhà máy rượu, cứ 10 tấn tinh bột (chứa 6,85\% tạp chất trơ) sẽ sản xuất được $7{,}21\text{ m}^3$ ethanol $40^\circ$ (cho khối lượng riêng của ethanol nguyên chất là $0{,}789\text{ g/cm}^3$). Hiệu suất của quá trình sản xuất là bao nhiêu phần trăm? \textit{(Kết quả làm tròn đến hàng đơn vị)}
\loigiai{
1. Khối lượng tinh bột nguyên chất:
\[m = 10 \times (100\% - 6{,}85\%) = 9{,}315\text{ tấn} = 9315\text{ kg}.\]
2. Thể tích ethanol nguyên chất:
\[V = 7{,}21\text{ m}^3 \times \dfrac{40}{100} = 2{,}884\text{ m}^3 = 2884\text{ lít} = 2\,884\,000\text{ cm}^3.\]
3. Khối lượng ethanol thực tế thu được:
\[m_{\text{TT}} = 2884 \times 0{,}789 = 2275{,}476\text{ kg}.\]
4. Theo sơ đồ phản ứng lý thuyết: $\ce{(C6H10O5)_n -> 2n C2H5OH}$
\[m_{\text{LT}} = 9315 \times \dfrac{2 \times 46}{162} = 5290\text{ kg}.\]
5. Hiệu suất của toàn bộ quá trình sản xuất:
\[H = \dfrac{m_{\text{TT}}}{m_{\text{LT}}} \times 100\% = \dfrac{2275{,}476}{5290} \times 100\% \approx 43{,}01\% \approx \mathbf{43\%}.\]
\textbf{Đáp số:} \textbf{43}.
}
\end{ex}

\begin{ex}
Một loại chất béo có chứa 80\% triolein về khối lượng. Xà phòng hóa hoàn toàn $22{,}1\text{ kg}$ chất béo này trong dung dịch $\ce{NaOH}$, đun nóng thu được $x$ bánh xà phòng. Biết rằng trong mỗi bánh xà phòng có chứa $60\text{ gam}$ sodium oleate. Xác định giá trị của $x$?
\begin{center}
	\includegraphics[width=4.5cm]{San_Pham/images_crop/fig_cau3_p3_soap.png}
\end{center}
\loigiai{
1. Khối lượng triolein trong $22{,}1\text{ kg}$ chất béo:
\[m_{\text{triolein}} = 22{,}1 \times 80\% = 17{,}68\text{ kg} = 17680\text{ gam}.\]
2. Số mol triolein: $n_{\text{triolein}} = \dfrac{17680}{884} = 20\text{ mol}$.\\
3. Phương trình xà phòng hóa:
\[\ce{(C17H33COO)3C3H5 + 3NaOH -> 3C17H33COONa + C3H5(OH)3}\]
\[\Rightarrow n_{\text{sodium oleate}} = 3 \times 20 = 60\text{ mol}.\]
4. Khối lượng sodium oleate: $m = 60 \times 304 = 18240\text{ gam}$.\\
5. Số bánh xà phòng thu được:
\[x = \dfrac{18240}{60} = \mathbf{304}\text{ (bánh)}.\]
\textbf{Đáp số:} \textbf{304}.
}
\end{ex}

\begin{ex}
Khi xà phòng hóa triglyceride X bằng dung dịch $\ce{NaOH}$ dư, đun nóng, thu được sản phẩm gồm glycerol, sodium stearate và sodium palmitate. Có bao nhiêu đồng phân cấu tạo thỏa mãn tính chất trên của X?
\loigiai{
Triglyceride X tạo thành từ 2 loại acid béo là stearic acid (gốc S) và palmitic acid (gốc P). Do đó trong phân tử X phải có cả gốc S và gốc P:
\begin{itemize}
	\item Dạng 2 gốc S + 1 gốc P:
	\begin{itemize}
		\item Gốc P ở vị trí trung tâm ($\beta$): S -- $\beta$(P) -- S (1 đồng phân).
		\item Gốc P ở vị trí đầu mạch ($\alpha$): P -- $\beta$(S) -- S (1 đồng phân).
	\end{itemize}
	\item Dạng 1 gốc S + 2 gốc P:
	\begin{itemize}
		\item Gốc S ở vị trí trung tâm ($\beta$): P -- $\beta$(S) -- P (1 đồng phân).
		\item Gốc S ở vị trí đầu mạch ($\alpha$): S -- $\beta$(P) -- P (1 đồng phân).
	\end{itemize}
\end{itemize}
Tổng số đồng phân cấu tạo thỏa mãn là $2 + 2 = \mathbf{4}$ đồng phân.\\
\textbf{Đáp số:} \textbf{4}.
}
\end{ex}

\vspace{0.4cm}
\noindent
\textbf{B. PHẦN TỰ LUẬN (3 CÂU -- 3 ĐIỂM)}

\setcounter{ex}{0}

\begin{ex}
\textbf{(1 điểm)} Viết phương trình phản ứng hoàn thành sơ đồ chuyển hóa sau (ghi rõ điều kiện phản ứng nếu có):
\begin{center}
\textbf{Tinh bột $\rightarrow$ Glucose $\rightarrow$ Ethanol $\rightarrow$ Acetic acid $\rightarrow$ Ethyl acetate.}
\end{center}
\loigiai{
Các phương trình hóa học:
\begin{enumerate}[(1)]
	\item $\ce{(C6H10O5)_n + n H2O} \xrightarrow{\text{H}^+, t^\circ} \ce{n C6H12O6}$
	\item $\ce{C6H12O6} \xrightarrow[\text{men rượu}]{30 - 35^\circ\text{C}} \ce{2 C2H5OH + 2 CO2 ^}$
	\item $\ce{C2H5OH + O2} \xrightarrow{\text{men giấm}} \ce{CH3COOH + H2O}$
	\item $\ce{CH3COOH + C2H5OH} \xrightleftharpoons[\text{H}_2\text{SO}_4\text{ đặc}, t^\circ]{} \ce{CH3COOC2H5 + H2O}$
\end{enumerate}
}
\end{ex}

\begin{ex}
\textbf{(1 điểm)} Aniline có nhiều ứng dụng quan trọng và đa dạng trong các ngành công nghiệp, đặc biệt là ngành sản xuất phẩm nhuộm và dược phẩm. Aniline có thể được tổng hợp từ benzene theo sơ đồ chuyển hoá sau:
\begin{center}
	\includegraphics[width=13.5cm]{San_Pham/images_crop/fig_cau2_tuluan_aniline.png}
\end{center}
Theo sơ đồ trên, từ 1 tấn benzene sẽ điều chế được bao nhiêu kg aniline? Biết hiệu suất toàn bộ quá trình là 60\%.
\loigiai{
Sơ đồ bảo toàn vòng benzene: $\ce{C6H6 -> C6H5NH2}$\\
Khối lượng mol: $M_{\ce{C6H6}} = 78\text{ g/mol}$; $M_{\ce{C6H5NH2}} = 93\text{ g/mol}$.\\
1. Khối lượng benzene: $m_{\text{benzene}} = 1\text{ tấn} = 1000\text{ kg}$.\\
2. Khối lượng aniline thu được theo lý thuyết (hiệu suất 100\%):
\[m_{\text{aniline (LT)}} = 1000 \times \dfrac{93}{78} \approx 1192{,}308\text{ kg}.\]
3. Do hiệu suất toàn bộ quá trình là 60\%, khối lượng aniline thực tế thu được là:
\[m_{\text{aniline (TT)}} = 1192{,}308 \times 60\% = \mathbf{715{,}38\text{ kg}} \text{ (hoặc } 715{,}4\text{ kg)}.\]
\textbf{Đáp số:} Khoảng $\mathbf{715{,}38\text{ kg}}$ aniline.
}
\end{ex}

\begin{ex}
\textbf{(1 điểm)}
\begin{enumerate}[a,]
	\item Trong cây xanh, glucose được tổng hợp nhờ quá trình quang hợp. Viết phương trình phản ứng xảy ra?
	\item Một nhóm học sinh muốn thử nghiệm phản ứng tráng bạc lên kính bằng nguyên liệu đầu là glucose. Giả sử lớp bạc có diện tích là $100\text{ cm}^2$ và độ dày là $0{,}5\text{ }\mu\text{m}$. Biết rằng khối lượng riêng của bạc là $10{,}49\text{ g/cm}^3$ và khối lượng mol của glucose là $180\text{ g/mol}$. Khối lượng glucose cần dùng là bao nhiêu gam? (Giả thiết hiệu suất phản ứng tráng bạc là 92\%).
\end{enumerate}
\loigiai{
a) Phương trình hóa học của quá trình quang hợp:
\[\ce{6 CO2 + 6 H2O} \xrightarrow[\text{diệp lục}]{\text{ánh sáng}} \ce{C6H12O6 + 6 O2 ^}\]
b) Khối lượng glucose cần dùng:
\begin{enumerate}[1.]
	\item Thể tích của lớp bạc cần mạ trên mặt kính:
	\[V = S \times h = 100\text{ cm}^2 \times (0{,}5 \cdot 10^{-4}\text{ cm}) = 0{,}005\text{ cm}^3.\]
	\item Khối lượng bạc cần tráng:
	\[m_{\ce{Ag}} = V \times D = 0{,}005\text{ cm}^3 \times 10{,}49\text{ g/cm}^3 = 0{,}05245\text{ gam}.\]
	\item Số mol bạc cần tạo thành: $n_{\ce{Ag}} = \dfrac{0{,}05245}{108}\text{ mol}$.
	\item Phương trình phản ứng tráng bạc của glucose: $\ce{C6H12O6 -> 2Ag}$
	\[\Rightarrow n_{\text{glucose (LT)}} = \dfrac{n_{\ce{Ag}}}{2} = \dfrac{0{,}05245}{2 \times 108}\text{ mol}.\]
	\item Khối lượng glucose theo lý thuyết:
	\[m_{\text{glucose (LT)}} = \dfrac{0{,}05245}{216} \times 180 = 0{,}043708\text{ gam}.\]
	\item Do hiệu suất phản ứng tráng bạc là 92\%, khối lượng glucose thực tế cần dùng là:
	\[m_{\text{glucose (cần dùng)}} = \dfrac{0{,}043708}{0{,}92} \approx \mathbf{0{,}0475\text{ gam}} \text{ (tương đương } 47{,}5\text{ mg)}.\]
\end{enumerate}
\textbf{Đáp số:} Khoảng $\mathbf{0{,}0475\text{ gam}}$ glucose.
}
\end{ex}

\begin{center}
\textbf{---------HẾT---------}
\end{center}

\end{document}
